Elemento opcional (Guia UECE 2026, 2.2.2.1.6), dirigido àqueles que contribuíram de forma relevante para a elaboração do trabalho: orientador(a), coorientador(a), membros da banca, professores, colegas, família, e agências ou instituições que financiaram ou viabilizaram a pesquisa (bolsas, laboratórios, coleta de dados etc.).

O título ``AGRADECIMENTOS'' é gerado automaticamente pelo macro \texttt{\textbackslash imprimiragradecimentos} (centralizado, caixa alta, negrito, sem pontuação); não repita esse título dentro do texto. O texto não deve ser dividido em parágrafos numerados ou alíneas -- normalmente uma frase por linha, uma para cada agradecimento, já é suficiente.

Regra que exige atenção do autor: o texto deve ser sucinto, claro e direto, evitando divagações, e não pode ultrapassar 2 (duas) laudas.

Substitua este parágrafo por frases como as do exemplo abaixo, adaptadas à sua realidade:

À Prof.ª Dr.ª [nome], pela orientação e confiança.

Aos membros da banca examinadora, pelas contribuições e correções sugeridas.

À [agência de fomento / instituição], pelo apoio recebido durante a pesquisa.

À minha família, pelo incentivo constante.
